% !TeX program = lualatex
% Vietnamese translation for OI Public Library
% Source-PDF: IOI中国国家候选队论文/2000/施遥论文.pdf
\documentclass[11pt]{article}
\usepackage{oipl-vn}

\title{Ứng dụng trí tuệ nhân tạo trong chương trình cờ vây}
\author{Shi Yao}
\date{}

\begin{document}
\maketitle

\origtitle{Application of artificial intelligence in Go programs}
\sourcepath{IOI China candidate team papers / 2000 / Shi Yao.pdf}

\section*{Từ khóa}

Cờ vây; tìm kiếm; khớp mẫu; cây trò chơi.

\section*{Tóm tắt}

Việc viết chương trình cờ vây được gọi là ``đá thử vàng'' của trí tuệ nhân
tạo, và là một bài toán khó lớn của kỹ thuật trí tuệ nhân tạo.

Bài viết giới thiệu ứng dụng và sự phát triển của trí tuệ nhân tạo trong
chương trình cờ vây; so sánh sự khác nhau và độ phức tạp giữa các thuật toán
trò chơi cho cờ vây và cờ vua; từ đó phân tích các điểm khó của thuật toán cờ
vây, đồng thời thảo luận cách phối hợp nhiều thuật toán trò chơi, gồm lý
thuyết điểm khí, khớp mẫu và cây trò chơi, trong chương trình cờ vây. Cuối
cùng bài viết đưa ra một ví dụ thuật toán cho chương trình xử lý sống chết
trong cờ vây để tham khảo.

\section{Khái quát}

\subsection{Giới thiệu cờ vây}

Tương truyền cờ vây do vua Nghiêu sáng tạo. Bàn cờ có mười chín đường ngang
dọc; thiên nguyên được xem như thái cực, thái cực sinh lưỡng nghi, tức quân
đen và quân trắng; lưỡng nghi sinh tứ tượng, tức bốn góc bàn cờ. Sách cổ
\textit{Dịch chỉ} của Ban Cố viết rằng cờ có trắng đen để phân âm dương, các
quân bày thành hàng lối như mô phỏng thiên văn. Từ đó có thể thấy cờ vây ban
đầu được tạo ra theo hình tượng bầu trời, rồi dần dần phát triển thành một
trò chơi đối kháng.

\subsection{Máy tính và cờ vây}

Việc dùng máy tính cho các trò chơi cờ gần như có lịch sử dài bằng chính máy
tính. Nội dung này chủ yếu thuộc kỹ thuật trí tuệ nhân tạo. Trí tuệ nhân tạo,
với tư cách một ngành khoa học, lần đầu được nêu ra vào thập niên 1950, và
ngay sau đó đã được áp dụng vào các trò chơi cờ.

Nhờ tiến bộ kỹ thuật, tốc độ máy tính tăng lên và thuật toán liên tục phát
triển, trình độ cờ vua máy tính hiện nay đã rất cao; trái lại, trình độ cờ
vây máy tính vẫn chững lại.

Với cờ vây, con người chơi cờ dựa vào kinh nghiệm, tức ``cảm giác cờ''. Ưu
thế của con người là phán đoán mơ hồ và trực giác nhạy bén; cao thủ thường có
thể chợt nảy ra một nước hay. Tất nhiên sự vật có hai mặt: cảm xúc và trực
giác của con người đôi khi cũng dẫn họ đến sai lầm, và tâm lý của kỳ thủ là
một yếu tố hết sức quan trọng. Trực giác vừa là vũ khí của con người, vừa có
thể là nguyên nhân thất bại khi đối đầu với người khác.

Ưu thế của máy tính là tốc độ tính toán cao; nhược điểm là không giỏi phán
đoán mơ hồ, không biết chọn điểm dựa trên kinh nghiệm, nên lượng tìm kiếm trở
nên quá lớn. Máy tính không bị cảm xúc chi phối, không bị trực giác làm lạc
hướng, vì vậy nó có thể phân biệt tương đối chính xác vùng lớn nhỏ, rộng hẹp,
và việc đó cũng không tốn nhiều thời gian. Nhưng máy tính rốt cuộc không có
cảm giác cờ; nó không biết nước nào tốt, nước nào xấu, mà chỉ có thể thử từng
bước một. Hoặc là nó tốn thời gian khổng lồ mà chưa chắc hữu ích, hoặc là xử
lý qua loa, và kết quả có thể đoán trước được.

\section{Độ phức tạp của cờ vây}

Cả thế cờ toàn cục lẫn bài toán sống chết trong cờ vây đều có thể quy về ba
đặc điểm phức tạp sau.

\subsection{Tính mơ hồ}

Tên gọi ``cờ vây'' xuất phát từ ý nghĩa vây đất. Nếu hai bên ngay từ đầu đều
đặt quân sát nhau, tốc độ đi cờ, tức tốc độ chiếm đất, sẽ cực kỳ chậm. Vì thế
từ bố cục, trung bàn cho đến quan tử lớn, hai bên thường chỉ dựng lên một
đường viền đại khái, thậm chí có lúc đường viền cũng chưa rõ. Thế lực đen
trắng khó phân định, hình thế đan xen răng lược. Điều này gây khó khăn rất
lớn cho xử lý bằng máy tính.

\subsection{Tính lặp lại}

Trong cờ vua, quân một khi bị ăn thì bị nhấc khỏi bàn cờ vĩnh viễn. Trên bàn
cờ vây, vị trí đã bị ăn vẫn có thể đặt quân lại, thậm chí có thể ăn ngược đối
phương. Như vậy độ khó của tìm kiếm tăng lên rất nhiều. Chẳng hạn có dạng
nước hy sinh quân rồi ăn lại, khiến một khoảng không có thể nhiều lần đổi
chủ. Vì vậy ``quân chết'' chưa hẳn đã chết, còn ``quân sống'' đôi khi cũng
đáng lo. Trong xử lý bằng máy tính, không thể đơn giản xác định sống chết của
một đám quân và ảnh hưởng của nó đến xung quanh.

\subsection{Tính linh hoạt}

Trong cờ vua, sự linh hoạt nhiều nhất thể hiện qua đổi quân. Câu ``thà bỏ một
quân chứ không mất tiên thủ'' cũng chỉ nói đến một quân; nếu là hai quân, ba
quân thì sao? Khi hai bên ngang sức, e rằng sẽ thua chắc. Hơn nữa, trong
thuật toán cờ vua thường có cách quy đổi các loại quân thành một giá trị nhất
định rồi cộng lại. Ví dụ trong cờ vua quốc tế, lấy một tốt làm đơn vị thì một
mã xấp xỉ ba tốt rưỡi, một xe xấp xỉ năm tốt rưỡi, v.v.; nhiều nhất lại điều
chỉnh thêm theo vị trí của quân.

Sự linh hoạt của cờ vây vượt xa mức đó. Có lúc bỏ hơn mười quân để lấy thế,
hoặc bỏ một góc hai, ba chục mục để chuyển đổi. Hơn nữa, giá trị của quân cờ
vây rất khó ước lượng; giá trị ấy không hoàn toàn nằm ở bản thân quân cờ, mà
thường phụ thuộc vào bố trí xung quanh, có ảnh hưởng thậm chí xiên qua cả bàn
cờ.

\section{Thuật toán trò chơi và so sánh giữa cờ vua với cờ vây}

Năm 1997, máy tính Deep Blue của IBM đánh bại Kasparov, gây chấn động thế
giới. Thực ra trình độ cờ vua máy tính đã chen chân vào hàng cao thủ thế giới
từ thập niên 1970 và 1980, còn phần mềm cờ tướng Trung Quốc cũng gần đạt
trình độ đại sư, vượt xa người chơi nghiệp dư thông thường. Riêng cờ vây vẫn
tiến rất chậm, đến mức thắng một kỳ thủ nghiệp dư yếu cũng chưa chắc, chưa
nói tới cao thủ hàng đầu. Nguyên nhân không ngoài sự sâu rộng, tinh vi và
biến hóa phức tạp của cờ vây: kỳ thủ dựa nhiều vào kinh nghiệm, trong khi máy
tính không có khả năng đó.

Hiện nay, thuật toán trò chơi cờ chủ yếu có hai loại lớn: khớp mẫu và dùng
cây trò chơi. Việc áp dụng chúng trong cờ vua có thể truy ngược đến thập niên
1950 và 1960, và đã rất thành công.

Cờ vây, cũng như cờ vua, là một trò chơi đối kháng. Nhìn bề ngoài chỉ có hai
loại quân đen trắng, tưởng như rất đơn giản; thực tế lại phức tạp hơn nhiều
so với các loại cờ có nhiều binh chủng.

Cờ vây máy tính bắt nguồn từ thập niên 1960. Hai tiến sĩ Zobrist và Ryder đều
đề cập đến chương trình cờ vây trong luận văn của mình: thuật toán của người
trước dựa trên nhận dạng mẫu, còn thuật toán của người sau dựa trên tìm kiếm,
tức dùng cây trò chơi. Hai loại thuật toán vốn có hiệu quả tốt trong cờ vua
lại thể hiện cực kém trong cờ vây, thậm chí không thắng nổi người chỉ có kinh
nghiệm vài ván.

Ta thử xét hai loại thuật toán trên.

Trước hết là khớp mẫu. Trong cờ vua, số quân ít và loại quân nhiều, nên việc
phân biệt và quy nạp tương đối dễ hơn, dù trên thực tế vẫn khó. Trong cờ vây,
rõ ràng không thể lưu trữ mọi mẫu. Vì thế chỉ còn cách khớp mơ hồ. Nhưng cách
này cũng có vẻ rất khó. Cờ vây có khái niệm \term{hình ngu}, thường chỉ hình
có hiệu suất kém. Thế nào là hiệu suất kém? Đó là lãng phí lực quân, đặt quân
ở nơi không cần thiết, dù chỉ lãng phí một quân. Khi khớp mơ hồ, chỉ khác một
quân thì phải phán đoán thế nào: là hình tốt hay hình ngu? Đây là một mâu
thuẫn. Huống hồ trong thực chiến còn có những ``diệu thủ trong hình ngu'', nên
càng khó phán đoán.

Tiếp theo là thuật toán tìm kiếm. Chi phí tìm kiếm cực lớn. Theo ước tính,
cờ vua khi tìm kiếm 7 hiệp có khoảng 50 đến 60 tỷ lựa chọn, tất nhiên là sau
khi đã tỉa cây trò chơi. Con số đó tuy rất lớn nhưng với trình độ kỹ thuật
hiện nay vẫn có thể chịu được. Còn trong cờ vây, chỉ cần tính sơ bộ: giả sử
mỗi nước chỉ có khoảng một trăm điểm khả thi, vốn đã rất ít, thì số biến hóa
trong 14 nước, tức 7 hiệp, xấp xỉ \(10^{28}\). Đây còn là ước lượng bảo thủ,
nhưng đã đủ đáng kinh ngạc. Có thể thấy dùng thuật toán tìm kiếm trên toàn
cục là không khả thi. Vì vậy cách làm hiện nay thường chỉ dùng cây trò chơi
để tìm kiếm trong các cục bộ có mục tiêu rõ ràng, như làm sống, giết quân,
phá vây, cắt đứt.

\section{Thuật toán cờ vây}

Hiện nay, các phần mềm cờ vây phổ biến trên thế giới chủ yếu được cấu thành
từ ba loại thuật toán sau.

\begin{enumerate}
  \item Cho mỗi quân cờ sinh ra một loại ảnh hưởng lên xung quanh. Ảnh hưởng
  này giảm khi khoảng cách tăng, rồi được tính và chồng lên nhau theo một
  công thức nhất định để phán đoán hình thế và ước lượng giá trị điểm đi. Điều
  này phù hợp với kỳ lý của cờ vây: với mỗi quân có thể ước lượng
  \term{thế lực} của nó. Trong nhóm này có lý thuyết \term{điểm khí} nổi
  tiếng.

  \item Xây dựng thư viện mẫu, lưu một lượng lớn mẫu như định thức và hình cờ
  để khớp. Việc này thực ra liên quan đến nhiều câu kỳ ngạn và kỳ lý trong cờ
  vây, chẳng hạn ``đầu hai quân phải bẻ'', ``bị trấn thì ứng bằng bay'',
  ``gặp cắt thì kéo dài một phía'', điểm vào trung tâm ba quân thẳng, điểm vào
  hình vuông, v.v. Tất cả đều được thiết kế theo tình huống cụ thể của cờ vây.

  \item Với các cục bộ có mục tiêu rõ ràng, dùng phương pháp tìm kiếm trong
  trí tuệ nhân tạo để thăm dò kết quả.
\end{enumerate}

Nói chung, hiện vẫn chưa tìm ra thuật toán có tính đột phá; ta chỉ có thể
tinh chỉnh kỹ trong ba loại thuật toán trên.

\section{Nhận dạng hình cờ vây}

Trong cờ vây, trình độ cao thấp phần lớn dựa vào cảm giác và kinh nghiệm của
kỳ thủ, mà cảm giác ấy chính là biểu hiện chủ yếu của trình độ. Cảm giác về
cờ thực ra là nhận thức chủ quan của con người đối với bản thân ván cờ, tức
hình cờ. Đây là một quá trình nhận dạng.

Năng lực nhận dạng là một đặc trưng lớn của trí tuệ. Từ thấp đến cao, năng
lực nhận dạng có thể chia thành ba tầng: tầng thiết bị là nhận dạng vật lý;
tầng động vật là nhận dạng mơ hồ; tầng con người là nhận dạng cảm xúc.

Nhận dạng vật lý là lượng hóa thông tin thu nhận được theo vật lý, hóa học và
sinh học. Việc này không cần kinh nghiệm hay trí tuệ, nên là mức nhận dạng
thấp nhất.

Nhận dạng mơ hồ là nhận ra phần hữu dụng trong một lượng lớn thông tin phức
tạp, tức liên hệ thông tin tiếp nhận với ký ức và kinh nghiệm trước đây, rồi
loại bỏ thông tin không liên quan.

Nhận dạng cảm xúc là nhận dạng bậc cao nhất. Nó hoàn toàn là nhận dạng cảm
tính.

Có thể nói trong cờ vây vừa có nhận dạng mơ hồ, vừa có nhận dạng cảm xúc. Với
một hình cờ mới, con người so sánh nó với các hình cờ trong kinh nghiệm, tổng
hợp tình hình xung quanh rồi đưa ra phán đoán. Quyết sách ấy mang khuynh
hướng cảm xúc cá nhân rõ rệt, liên quan đến cách hiểu của từng người, và rất
nhiều trường hợp không có kết luận dứt khoát. Vì vậy ngay cả ván cờ giữa các
cao thủ cửu đẳng cũng có thể hoàn toàn khác nhau, khó nói bên nào hơn bên
nào.

Đối với máy tính, việc phán đoán mơ hồ hình cờ lại rất khó, vì đây không phải
âm thanh cũng không phải hình ảnh theo nghĩa thông thường; chỉ khác một quân,
tốt xấu đã có thể cách biệt rất xa.

Nói chung, hiện nay máy tính chỉ có thể nhận dạng đơn giản hình cờ vây để
giảm tìm kiếm. Thực chất đây là cách trao cho máy tính một mức ``kinh
nghiệm'' nhất định.

\section{Thuật toán và cài đặt cho bài toán sống chết trong cờ vây}

Sống chết là một bài toán điển hình trong cờ vây, cũng có thể xem là một mô
hình thu nhỏ của thuật toán cờ vây. Nó cần phối hợp cả ba loại thuật toán đã
nêu. Hiện nay phần mềm sống chết đã đạt trình độ khá cao, chủ yếu vì đây chỉ
là một bài toán cục bộ; còn quan hệ chằng chịt với toàn cục thì vẫn cực kỳ
khó nắm bắt.

\paragraph{Chú thích.}
Để trình bày đơn giản, trong phần dưới đây đen luôn là bên tấn công, trắng
luôn là bên muốn làm sống; phần này chỉ lấy trắng làm ví dụ thiết kế thuật
toán. Các thuật toán đều đã được giản lược, mục đích chỉ là minh họa phương
pháp. Do thời gian và trình độ có hạn, thuật toán trong phần này khó tránh
thiếu sót.

Với bài toán sống chết trong cờ vây, trước hết hãy xem cách con người suy
nghĩ khi chơi cờ.

Trước tiên xem có bao nhiêu vị trí mắt và đã sống hay chưa. Nếu chưa sống,
tính còn thiếu bao nhiêu mắt và có thể bổ sung được không. Trong quá trình
làm mắt hoặc phá mắt, con người luôn dùng cảm giác đầu tiên làm manh mối để
tính toán. Cảm giác đầu tiên ấy từ đâu mà ra? Thực ra nó là kinh nghiệm dài
lâu trong nhận thức hình cờ: ai nhiều kinh nghiệm thì mạnh, ai thiếu kinh
nghiệm thì yếu. Trong sống chết có nhiều hình cờ thường gặp, và người xưa đã
tổng kết một số kỳ ngạn như ``giết cờ dùng bẻ'', ``điểm hai một thường có
diệu thủ''; ngoài ra những thủ pháp như kẹp, điểm, nhảy qua cũng thường tạo
thành nước giết. Một số điểm đi thì không cần xét. Vì sống chết là một bài
toán cục bộ có mục tiêu rõ ràng, tìm kiếm là chủ đạo. Do đó chất lượng của
một phần mềm sống chết cờ vây nằm chủ yếu ở tối ưu hóa và tỉa nhánh trong tìm
kiếm.

Thuật toán của tôi cho bài toán sống chết gồm hai phần lớn:
\begin{itemize}
  \item phán đoán tĩnh vị trí mắt;
  \item tìm kiếm.
\end{itemize}

Tuyến chính của chương trình là tìm kiếm, nhưng phán đoán tĩnh vị trí mắt
được xen vào trong quá trình đó.

\subsection{Phán đoán tĩnh vị trí mắt}

\subsubsection{Phân chia thế lực}

Như đã nói ở trên, ranh giới giữa đen và trắng trong cờ vây là mơ hồ, rất
khó phân chia chính xác. Nếu không thể phân chia chính xác thì máy tính khó
xử lý. Vì vậy ta dùng phương pháp của lý thuyết điểm khí để xấp xỉ tính
\term{thế lực}.

Phân chia thế lực xuyên suốt toàn bộ phán đoán tĩnh vị trí mắt và có vai trò
rất lớn.

Một quân cờ có ảnh hưởng nhất định đến xung quanh, gọi là \term{thế lực};
các vị trí xung quanh một quân gọi là \term{điểm khí}. Trong hình minh họa
gốc, các vị trí gần quân được chia theo các mức sau:
\begin{enumerate}
  \item vị trí kề cạnh là điểm khí cấp \(1\);
  \item vị trí chéo nhọn là điểm khí cấp \(1.5\), còn vị trí bẻ là điểm khí
  cấp \(2\);
  \item vị trí cách một đường thẳng là điểm khí cấp \(2\);
  \item vị trí bay nhỏ là điểm khí cấp \(2.5\).
\end{enumerate}

Ảnh hưởng ứng với từng cấp điểm khí được cho trong bảng sau.
\[
\begin{array}{c|cccc}
\text{Cấp điểm khí} & 1 & 1.5 & 2 & 2.5\\
\hline
\text{Trị tuyệt đối của ảnh hưởng} & 6 & 5 & 3 & 2
\end{array}
\]
Ảnh hưởng của quân đen lấy giá trị dương, ảnh hưởng của quân trắng lấy giá
trị âm.

Dĩ nhiên, giữa các vị trí nói trên và vị trí ban đầu không được có quân cờ
ngăn cách; nếu có, bất kể đen hay trắng, thế lực đều bị chặn. Ngoài ra cần bổ
sung rằng ảnh hưởng của một màu lên một điểm chỉ lấy từ quân cùng màu gần
nhất, tức quân có ảnh hưởng lớn nhất.

Như vậy, một vùng nửa mở cũng có thể bị thế lực khép lại. Một số quân nằm
trong vòng vây dày đặc, vì thế lực đối phương xung quanh quá mạnh, giá trị
thế lực của chúng thậm chí bị đối phương đảo ngược và tự nhiên trở thành
quân chết. Điều này không tuyệt đối; dạng hy sinh rồi ăn lại là ví dụ quân
chết có thể sống lại.

Bài viết này chỉ trình bày một phương pháp cơ bản. Khi thao tác thật có thể
biến đổi tùy tình huống; với phần mềm cờ vây tốt, ảnh hưởng cũng chưa chắc
được chồng tuyến tính.

\subsubsection{Xác định vị trí mắt của trắng}

Trước hết, ta biết mỗi đám quân cần ít nhất hai mắt thật để sống. Sự khác
nhau giữa mắt thật và mắt giả được minh họa trong PDF gốc bằng ba thế ở góc,
biên và trung tâm. Diễn giải bằng chữ như sau: ở góc, điểm chéo \(A\) phải
do bên mình chiếm thì mắt mới là mắt thật, ngược lại là mắt giả; ở biên, hai
điểm chéo \(A,B\) đều phải do bên mình chiếm thì mắt mới là mắt thật; ở trung
tâm, trong bốn điểm chéo \(A,B,C,D\), bên mình phải chiếm ít nhất ba điểm thì
mới thành mắt thật, ngược lại là mắt giả. Nếu bên mình đã chiếm hai điểm, còn
hai điểm kia trống, đó cũng là mắt thật, vì hai điểm này là miai: chắc chắn
sẽ lấy được một trong hai.

Bây giờ xét thế nào là ``chiếm''. Hàm nghĩa của ``chiếm'' đối với đen và
trắng khác nhau.

Đối với trắng, chiếm không nhất thiết nghĩa là tại điểm đó có quân trắng; chỉ
cần thế lực trắng tại điểm đó cực mạnh. Khi đó dù đen đặt quân vào đây, quân
đen cũng sẽ bị trắng ăn, cuối cùng điểm đó vẫn thuộc về trắng.

Đối với đen, chiếm nghĩa là điểm đó có quân đen. Ở nơi không có quân đen, dù
thế lực mạnh đến đâu, một khi trắng đặt quân vào đó, đen cũng không có cách
nào ngay lập tức xử lý.

Thứ hai, cần nói rõ rằng vị trí mắt không nhất thiết chỉ có hai trạng thái
thật hoặc giả. Hình minh họa trong PDF gốc cho thấy: nếu trắng đi trước tại
điểm \(A\) thì có thể bổ thành mắt thật, ngược lại nếu đen đi trước thì thành
mắt giả. Ta tạm gọi đó là \term{nửa mắt}.

\paragraph{Phán đoán mắt một mục.}
Mắt một mục tương đối dễ phán đoán: chỉ cần tính như trên, chia thành ba
loại là mắt thật, mắt giả không cần tính, và mắt thật khi trắng đi trước.

\paragraph{Phán đoán khoảng không từ hai đến năm mục.}
Vị trí mắt từ hai mục trở lên có hình dạng phức tạp, khó phán đoán đơn giản.
Tuy hình dạng phức tạp, số lượng hình cơ bản vẫn hữu hạn: khoảng hơn mười
hình, sau xoay và lật thành ba, bốn chục hình. Vì vậy có thể làm thư viện
hình cờ để khớp. Có năm loại: trắng đi trước được một mắt thật; luôn là một
mắt thật; trắng đi trước được hai hoặc hơn hai mắt thật; luôn là hai hoặc hơn
hai mắt thật; và không có mắt thật, loại này có thể bỏ qua.

\paragraph{Phán đoán khoảng không trên năm mục.}
Khoảng không trên năm mục có rất nhiều hình dạng, không thể kể hết, nên rất
khó làm thành thư viện mẫu. Nhưng điều đáng chú ý là khoảng không trên năm
mục, nếu hình dạng hoàn chỉnh, gần như đều có thể thành hai mắt thật. Vì vậy
khi phán đoán, tạm thời có thể xem nó là hai hoặc hơn hai mắt thật.

\subsubsection{Phán đoán sống chết}

Sau khi xác định được vị trí mắt như trên, ta có thể đi vào điểm mấu chốt.
Mục đích chính của phán đoán tĩnh vị trí mắt là xác định loại sống chết:
trắng đi trước thì sống, trắng đi trước vẫn chết, và trắng đi sau vẫn sống.

\paragraph{Trắng đi trước thì sống.}
Trắng đi trước có thể làm sống, đi sau thì chết. Điều kiện là một trong các
trường hợp sau:
\begin{itemize}
  \item một mắt thật và một nửa mắt;
  \item ba nửa mắt.
\end{itemize}

\paragraph{Trắng đi trước vẫn chết.}
Ngay cả khi trắng đi trước cũng chết. Điều kiện là một trong các trường hợp
sau:
\begin{itemize}
  \item một hoặc không có mắt thật, và không có nửa mắt;
  \item không có mắt thật, có hai hoặc ít hơn hai nửa mắt.
\end{itemize}

\paragraph{Trắng đi sau vẫn sống.}
Ngay cả khi trắng đi sau cũng là cờ sống. Điều kiện là một trong các trường
hợp sau:
\begin{itemize}
  \item hai hoặc hơn hai mắt thật;
  \item một mắt thật và hai hoặc hơn hai nửa mắt;
  \item bốn hoặc hơn bốn nửa mắt.
\end{itemize}

Nói đơn giản, lấy một mắt thật quy đổi thành hai nửa mắt. Nếu ít hơn ba nửa
mắt thì trắng đi trước vẫn chết; đúng bằng ba nửa mắt thì trắng đi trước mới
sống; lớn hơn ba nửa mắt thì trắng đi sau vẫn sống.

\subsection{Tìm kiếm bằng cây trò chơi}

Thuật toán tìm kiếm là tuyến chính, nhưng không phức tạp bằng phán đoán tĩnh
vị trí mắt.

Trong cây trò chơi, các tầng lẻ là lượt trắng, tức máy tính, còn các tầng
chẵn là lượt đen. Sơ đồ cây chỉ cần biểu diễn các lượt trắng--đen luân phiên;
trong cài đặt, với mỗi nút ta dùng khả năng trắng sống trong trạng thái đó làm
hàm đánh giá.

Với việc chọn điểm đi cho trắng ở tầng lẻ, có thể đặt các hạn chế sau:
\begin{enumerate}
  \item không đi ra ngoài vòng vây của đen;
  \item quân này không được bị đen ăn bằng một nước, trừ khi cả đám quân bị
  ăn có ít hơn bốn quân, vì nếu bị ăn từ bốn quân trở lên thì đó có thể là
  diệu thủ hy sinh rồi ăn lại.
\end{enumerate}

Với việc chọn điểm đi cho đen ở tầng chẵn, giới hạn trong vòng vây của đen và
bên trong vòng vây ấy.

Sau khi liệt kê các điểm đi của trắng, mỗi điểm sinh ra một trạng thái. Nếu
trạng thái này là ``trắng đi sau vẫn sống'', lập tức trả về trắng sống ở
trạng thái đó, tức khả năng trắng sống bằng \(1\). Nếu không, tìm xem trong
trạng thái đó có nước đen nào khiến trắng rơi vào loại ``trắng đi trước vẫn
chết'' hay không. Nếu có, nút này không cần mở rộng tiếp và trả về trắng chết
ở trạng thái đó, tức khả năng trắng sống bằng \(0\). Nếu không thuộc hai tình
huống trên thì mở rộng tầng tiếp theo, tức tầng đen, rồi dùng độ lớn khả năng
trắng sống trong các trạng thái tầng dưới để đánh giá nút này.

Sau khi liệt kê các điểm đi của đen, với mỗi điểm đi của đen cũng sinh ra
một trạng thái. Nếu trong trạng thái đó trắng thuộc loại ``trắng đi trước thì
sống'', lập tức trả về trắng sống, tức khả năng trắng sống bằng \(1\). Nếu
không, tìm xem liệu trắng có nước nào khiến trạng thái trở thành ``trắng đi
sau vẫn sống'' hay không. Nếu có, nút này không cần mở rộng tiếp và trả về
trắng sống, tức khả năng trắng sống bằng \(1\). Nếu không thuộc các tình
huống trên thì tiếp tục mở rộng nút, rồi dùng độ lớn khả năng trắng chết
trong các trạng thái tầng dưới để đánh giá nút này.

Cuối cùng, chương trình căn cứ vào khả năng trắng sống và chọn điểm đi có giá
trị lớn nhất. Tất nhiên trong cây trò chơi nên dùng tỉa \(\alpha\)-\(\beta\):
ở tầng trắng, loại bỏ các nút có khả năng trắng sống nhỏ và không mở rộng;
ở tầng đen, loại bỏ các nút có khả năng trắng sống lớn và không mở rộng. Một
số vấn đề cụ thể về cây trò chơi có thể tham khảo tài liệu liên quan.

\section{Triển vọng}

Ngoài các phương pháp trên, trong lĩnh vực trí tuệ nhân tạo còn có một số
phương án khả thi sau.

\paragraph{Mạng nơ-ron nhân tạo.}
Sự phát triển của mạng nơ-ron nhân tạo trong những năm gần đây đã mở thêm một
con đường cho chương trình cờ vây: thêm phần mạng nơ-ron nhân tạo để phần mềm
có năng lực học tập, ghi nhớ và liên tưởng nhất định; cho nó đối cục với cao
thủ để hấp thu một lượng ``tri thức'' và nuôi dưỡng cảm giác cờ. MFGO của
David Fotland tại Hoa Kỳ đã có một số thử nghiệm hữu ích theo hướng này,
nhưng hiện có thể vẫn chưa hoàn thiện. David Fotland, kỹ sư của Hewlett
Packard và là người thiết kế phần mềm Many Faces of Go, từng nói đại ý rằng
tìm kiếm brute force hoàn toàn không có tác dụng với cờ vây; phải tạo ra một
chương trình tinh khôn như con người. Quả thật, cờ vây phức tạp đến mức chỉ
dạy máy tính chơi cờ một cách máy móc là hoàn toàn không đủ; nên dùng mạng
nơ-ron nhân tạo để làm máy tính ``thông minh'' hơn, để nó chủ động học tập.
Nhưng kỹ thuật ở hướng này chưa đạt đến mức cao, nên hiện vẫn chỉ là một ý
tưởng.

\paragraph{Hệ chuyên gia.}
Hệ chuyên gia đã được dùng rộng rãi trong nhiều lĩnh vực. Thực ra trong cờ
vây cũng có thể xây dựng hệ chuyên gia. Có thể lập kho tri thức hình cờ, dựa
trên kỳ ngạn và các kinh nghiệm khác do con người tổng kết, rồi trao chúng
cho máy tính. Đây là một phương án tương đối hệ thống.

\paragraph{Hệ sinh.}
Với định thức và một số biến hóa thường gặp, có thể xây dựng hệ sinh để máy
tính tự sinh ra biến hóa. Điểm then chốt là chọn phương pháp điều khiển tốt
để phán đoán hay dở. Đây là một phương án tiền xử lý cho các bài toán cục bộ.

Cờ không có điểm tận cùng; trên chặng đường dài của cờ vây máy tính, điều đó
lại càng đúng.

Hiện nay, phần mềm cờ vây đã là một trong những mũi nhọn của kỹ thuật trí tuệ
nhân tạo, được gọi là ``đá thử vàng'' của trí tuệ nhân tạo. Muốn tiến thêm
một bước, không ngoài tối ưu hóa cả phần mềm lẫn phần cứng. Tối ưu hóa thuật
toán ngày càng đi vào chi tiết, còn đột phá căn bản vẫn chưa tìm thấy; tốc độ
phần cứng tuy tăng ngày càng nhanh, nhưng vẫn còn xa mới đáp ứng được yêu
cầu. HandTalk hiện là một trong các phần mềm cờ vây mạnh nhất thế giới, nhưng
theo tác giả của nó, giáo sư Chen Zhixing, nó chỉ có kỳ lực cấp 9 kyu. Giáo
sư Chen nói điểm nghẽn then chốt trong phát triển phần mềm cờ vây là ``cờ vây
quá phức tạp''.

Tuy trong bài toán sống chết, máy tính đã đạt đến một trình độ nhất định,
nhưng bản thân cờ vây là một học vấn lớn về sự chuyển hóa và cân bằng liên
tục giữa được và mất. Mất ở chỗ này nhưng được ở chỗ khác; quan hệ tinh vi
giữa được và mất hiện là điều máy tính khó nắm bắt. Tại giải cờ vây máy tính
cúp Ing, phóng viên hỏi các thí sinh còn bao lâu nữa máy tính mới thắng được
cao thủ con người. Câu trả lời là một trăm năm sau. Dù câu trả lời ấy bảo thủ
hay phóng đại, con đường của cờ vây máy tính vẫn còn rất dài. Khi một ngày
nào đó thật sự hoàn thành ước mơ của ông Ing Chang-ki rằng máy tính đánh bại
được cửu đẳng con người, đó mới là dấu hiệu máy tính thật sự đạt đến trí tuệ.

\section*{Tài liệu tham khảo}

\begin{enumerate}
  \item Liu Fusheng, Wang Jiande, \textit{Hướng dẫn Olympic Tin học Quốc tế
  dành cho thanh thiếu niên: tìm kiếm trí tuệ nhân tạo và thiết kế chương
  trình}.
  \item E. Rich, \textit{Dẫn nhập trí tuệ nhân tạo}.
  \item Chen Zhixing, \url{http://www.wulu.com}.
  \item Bản điện tử của \textit{Báo cờ vây}, \url{http://www.wqb.com.cn}.
  \item American Go Association, \url{http://www.usgo.org/computer}.
\end{enumerate}

\paragraph{Phụ lục.}
Chương trình minh họa sống chết trong cờ vây được nhắc tới trong nguồn gốc
nhưng không xuất hiện trong phần văn bản đọc được của PDF này.

\end{document}
